\section{Tính toán với kim loại và dung dịch}

\subsection{Nhắc lại bốn bước giải}

Bài toán kim loại vẫn dùng đúng cách giải chung đã học:
\begin{ghinho}{Các bước giải}
\begin{enumerate}
\item Tính số mol.
\item Viết phương trình hoặc sơ đồ phản ứng.
\item Tính các giá trị trung gian cần thiết: số mol, khối lượng, chất hết -- dư, ẩn số,...
\item Tính đại lượng đề hỏi và kết luận.
\end{enumerate}
\end{ghinho}

\subsection{Một kim loại -- một phản ứng}

\begin{vidu}
Cho \SI{6.50}{\gram} Zn tác dụng hoàn toàn với dung dịch \ce{HCl} dư. Tính thể tích \ce{H2} thu được ở điều kiện có $V_m=\SI{24.8}{\liter\per\mole}$.
\end{vidu}
\begin{loigiai}
\textbf{Bước 1. Tính số mol}
\[
n_{\ce{Zn}}=\frac{6.50}{65}=\SI{0.100}{\mole}.
\]

\textbf{Bước 2. Viết phương trình hóa học}
\[
\begin{array}{c@{\;}c@{\;}c@{\;}c@{\;}c}
\ce{Zn} & + & \ce{2HCl} & \ce{->} & \ce{ZnCl2 + H2}\\
0.100 & & & & 0.100
\end{array}
\]

\textbf{Bước 3. Tính giá trị trung gian}
\[
n_{\ce{H2}}=\SI{0.100}{\mole}.
\]

\textbf{Bước 4. Tính đại lượng đề hỏi}
\[
V_{\ce{H2}}=0.100\times24.8=\SI{2.48}{\liter}.
\]

\textbf{Vậy,} thể tích \ce{H2} thu được là \SI{2.48}{\liter}.
\end{loigiai}

\subsection{Chất hết và chất dư}

Khi cả hai chất phản ứng đều có lượng xác định, phải so sánh theo tỉ lệ hệ số trong phương trình trước khi tính sản phẩm.

\begin{vidu}
Cho \SI{6.50}{\gram} Zn vào \SI{100}{\milli\liter} dung dịch \ce{HCl} \SI{1.00}{\mole\per\liter}. Xác định chất hết, chất dư và thể tích \ce{H2} thu được.
\end{vidu}
\begin{loigiai}
\textbf{Bước 1. Tính số mol}
\[
n_{\ce{Zn}}=\SI{0.100}{\mole},\qquad n_{\ce{HCl}}=0.100\times1.00=\SI{0.100}{\mole}.
\]

\textbf{Bước 2. Viết phương trình hóa học}
\[
\begin{array}{c@{\;}c@{\;}c@{\;}c@{\;}c}
\ce{Zn} & + & \ce{2HCl} & \ce{->} & \ce{ZnCl2 + H2}\\
x & & 2x & & x
\end{array}
\]

\textbf{Bước 3. Tính các giá trị trung gian}
Theo phương trình, \SI{0.100}{\mole} HCl chỉ phản ứng với \SI{0.050}{\mole} Zn. Vì vậy HCl hết, Zn dư:
\[
n_{\ce{Zn},\text{dư}}=0.100-0.050=\SI{0.050}{\mole}.
\]
\[
n_{\ce{H2}}=\SI{0.050}{\mole}.
\]

\textbf{Bước 4. Tính đại lượng đề hỏi}
\[
V_{\ce{H2}}=0.050\times24.8=\SI{1.24}{\liter}.
\]

\textbf{Vậy,} HCl hết, Zn dư \SI{0.050}{\mole} và thu được \SI{1.24}{\liter} \ce{H2}.
\end{loigiai}

\subsection{Hỗn hợp hai kim loại -- lập hệ hai ẩn}

Khi hỗn hợp có hai thành phần chưa biết, thường đặt số mol là $x,y$. Một phương trình thường đến từ khối lượng hỗn hợp; phương trình còn lại đến từ khí, chất rắn, lượng acid hoặc một dữ kiện độc lập khác.

\begin{vidu}
Hòa tan hoàn toàn \SI{9.00}{\gram} hỗn hợp Mg và Al bằng \ce{HCl} dư, thu được \SI{11.16}{\liter} \ce{H2} ở điều kiện có $V_m=\SI{24.8}{\liter\per\mole}$. Tính khối lượng mỗi kim loại.
\end{vidu}
\begin{loigiai}
Gọi số mol Mg và Al lần lượt là $x,y$.

\textbf{Bước 1. Tính số mol}
\[
n_{\ce{H2}}=\frac{11.16}{24.8}=\SI{0.450}{\mole}.
\]

\textbf{Bước 2. Viết phương trình hóa học}
\[
\begin{array}{c@{\;}c@{\;}c@{\;}c@{\;}c}
\ce{Mg} & + & \ce{2HCl} & \ce{->} & \ce{MgCl2 + H2}\\
x & & & & x
\end{array}
\]
\[
\begin{array}{c@{\;}c@{\;}c@{\;}c@{\;}c}
\ce{2Al} & + & \ce{6HCl} & \ce{->} & \ce{2AlCl3 + 3H2}\\
y & & & & \frac{3y}{2}
\end{array}
\]

\textbf{Bước 3. Lập hệ và tìm ẩn}
\[
\begin{cases}
24x+27y=9.00,\\
x+\dfrac{3y}{2}=0.450.
\end{cases}
\Rightarrow
\begin{cases}
x=0.150,\\
y=0.200.
\end{cases}
\]

\textbf{Bước 4. Tính đại lượng đề hỏi}
\[
m_{\ce{Mg}}=0.150\times24=\SI{3.60}{\gram},
\]
\[
m_{\ce{Al}}=0.200\times27=\SI{5.40}{\gram}.
\]

\textbf{Vậy,} hỗn hợp chứa \SI{3.60}{\gram} Mg và \SI{5.40}{\gram} Al.
\end{loigiai}
